\documentclass[11pt,a4paper]{article}

% ==================================================
% PAQUETES
% ==================================================
\usepackage[utf8]{inputenc}
\usepackage[T1]{fontenc}
\usepackage[spanish]{babel}
\usepackage[margin=2.3cm]{geometry}
\usepackage{graphicx}
\usepackage{float}
\usepackage{hyperref}
\usepackage{xcolor}
\usepackage{booktabs}
\usepackage{longtable}
\usepackage{tabularx}
\usepackage{array}
\usepackage{amsmath}
\usepackage{enumitem}
\usepackage{pdflscape}
\usepackage{tikz}
\usepackage{caption}
\usepackage{fancyhdr}

\usetikzlibrary{
    arrows.meta,
    positioning,
    shapes.geometric,
    fit,
    backgrounds
}

\hypersetup{
    colorlinks=true,
    linkcolor=blue!50!black,
    citecolor=blue!50!black,
    urlcolor=blue!60!black,
    pdftitle={Framework y Aplicacion del Proceso de Requisitos},
}

\setlength{\parskip}{5pt}
\setlength{\parindent}{0pt}
\renewcommand{\arraystretch}{1.3}

\pagestyle{fancy}
\fancyhf{}
\fancyhead[L]{Ingenieria de Requisitos}
\fancyhead[R]{Almacenes Tia}
\fancyfoot[C]{\thepage}

\newcolumntype{Y}{>{\raggedright\arraybackslash}X}

% ==================================================
% DOCUMENTO
% ==================================================

\title{
    \textbf{Framework y Aplicacion del Proceso de Requisitos}\\
    \large Plataforma de Comercio Electronico de Almacenes Tia
}

\begin{document}
\newpage


% ==================================================
\section{Framework de Ingenieria de Requisitos segun Pohl}
% ==================================================

\subsection{Descripcion del framework}

De acuerdo con el framework de ingenieria de
requisitos propuesto por Klaus Pohl, la ingenieria
de requisitos debe abordar las necesidades de los
stakeholders mediante actividades sistematicas
que permitan obtener requisitos comprensibles,
consistentes y verificables.

El framework considera tres dimensiones
fundamentales:

\begin{enumerate}
    \item \textbf{Especificacion:} permite definir
    y documentar las necesidades y requisitos
    del sistema.

    \item \textbf{Representacion:} establece las
    formas mediante las cuales los requisitos
    pueden expresarse y comprenderse, utilizando
    descripciones textuales, modelos y diagramas.

    \item \textbf{Acuerdo:} busca establecer un
    entendimiento compartido entre los diferentes
    stakeholders respecto a los requisitos
    establecidos.
\end{enumerate}

Estas dimensiones se relacionan entre si y permiten
organizar el proceso de ingenieria de requisitos
durante el desarrollo de un sistema.

\subsection{Actividades principales (Core)}

Para la aplicacion realizada se consideran las
siguientes actividades principales:

\begin{itemize}

    \item \textbf{Elicitacion:}

    Consiste en descubrir y recopilar las necesidades,
    expectativas y restricciones de los stakeholders.
    En Almacenes Tia se realiza mediante la observacion
    de la interfaz publica, identificando las
    funcionalidades disponibles y las interacciones
    del usuario.

    \item \textbf{Analisis:}

    Consiste en examinar los requisitos identificados,
    detectar ambiguedades, conflictos, dependencias
    y establecer prioridades. En el sistema seleccionado
    se analiza la relacion entre la busqueda de productos,
    el catalogo y el carrito de compras.

    \item \textbf{Especificacion:}

    Consiste en documentar formalmente los requisitos
    mediante identificadores, descripciones, tipos,
    prioridades y criterios de aceptacion. Para ello
    se utilizan estructuras basadas en IEEE 29148
    y la sintaxis EARS.

    \item \textbf{Validacion:}

    Consiste en comprobar que los requisitos sean
    claros, completos, consistentes, verificables
    y que correspondan con las necesidades identificadas.
    Se utiliza una lista de verificacion y se contrastan
    los requisitos con las evidencias disponibles.

\end{itemize}

\subsection{Actividades transversales}

Las actividades transversales permiten mantener
la calidad y consistencia de los requisitos durante
todo el proceso.

\begin{itemize}

    \item \textbf{Gestion de requisitos:}

    Permite controlar identificadores, versiones,
    modificaciones y trazabilidad. En Almacenes Tia
    se aplica mediante una matriz que relaciona
    cada requisito con su fuente y evidencia.

    \item \textbf{Gestion de conflictos:}

    Permite identificar y resolver discrepancias
    entre las necesidades de los usuarios y las
    restricciones del sistema. Por ejemplo, puede
    existir una diferencia entre la expectativa
    de consultar productos disponibles inmediatamente
    y la informacion que realmente proporciona
    la interfaz.

\end{itemize}

\subsection{Relacion del framework con Almacenes Tia}

\begin{table}[H]
\centering
\caption{Aplicacion de las actividades del framework}
\small
\begin{tabularx}{\textwidth}{p{3cm}YY}
\toprule
\textbf{Actividad} &
\textbf{Aplicacion en Tia} &
\textbf{Resultado esperado} \\
\midrule

Elicitacion &
Exploracion del catalogo, buscador y carrito. &
Necesidades identificadas y evidencias. \\

Analisis &
Identificacion de ambiguedades, prioridades y
dependencias entre funcionalidades. &
Requisitos organizados y clasificados. \\

Especificacion &
Documentacion de requisitos funcionales y no
funcionales. &
Matriz de requisitos. \\

Validacion &
Revision de claridad, consistencia y verificabilidad. &
Requisitos revisados. \\

Gestion de requisitos &
Control de identificadores, fuentes y modificaciones. &
Trazabilidad documental. \\

Gestion de conflictos &
Analisis de diferencias entre expectativas y
restricciones observables. &
Conflictos identificados y alternativas propuestas. \\

\bottomrule
\end{tabularx}
\end{table}

% ==================================================
\section{Identificacion de actores y stakeholders}
% ==================================================

Los stakeholders representan a las personas o entidades
que tienen interes o influencia sobre el sistema.

A partir del alcance de la observacion se identifican
los siguientes actores y participantes:

\begin{table}[H]
\centering
\caption{Identificacion de actores}
\small
\begin{tabularx}{\textwidth}{p{3cm}YY}
\toprule
\textbf{Actor} &
\textbf{Responsabilidad o interes} &
\textbf{Necesidad principal} \\
\midrule

Cliente &
Utiliza la plataforma para consultar productos
y realizar compras. &
Encontrar productos y gestionar su compra. \\

Personal de tienda &
Participa en la atencion y gestion de productos
en los establecimientos. &
Contar con informacion adecuada sobre productos
y disponibilidad. \\

Administrador de la plataforma &
Gestiona aspectos operativos y funcionales
del sistema. &
Mantener el funcionamiento y la informacion
de la plataforma. \\

Empresa Almacenes Tia &
Organizacion responsable del servicio. &
Ofrecer un canal digital de comercio electronico. \\

\bottomrule
\end{tabularx}
\end{table}

\textbf{Nota:} los actores se identifican considerando
el contexto funcional del sistema. La participacion
de personal interno y administradores no fue
entrevistada directamente, por lo que sus necesidades
constituyen una identificacion preliminar.

% ==================================================
\section{Elicitacion de requisitos}
% ==================================================

\subsection{Tecnicas seleccionadas}

Para obtener informacion sobre el sistema se
consideraron las siguientes tecnicas:

\begin{enumerate}

    \item \textbf{Observacion participante:}

    Consiste en explorar directamente la interfaz
    publica e identificar las acciones disponibles
    para el usuario.

    \item \textbf{Exploracion analitica:}

    Consiste en revisar las funcionalidades y
    establecer las relaciones entre las diferentes
    operaciones de la plataforma.

    \item \textbf{Analisis documental:}

    Consiste en revisar la informacion publica
    disponible en el sitio web para comprender
    el contexto y funcionamiento general del servicio.

\end{enumerate}

\subsection{Procedimiento de observacion}

La observacion se plantea sobre el sitio web oficial
de Almacenes Tia, considerando el contexto local
de San Carlos, Quevedo.

Las actividades de exploracion comprenden:

\begin{enumerate}
    \item Acceder al sitio web oficial.
    \item Examinar la presentacion del catalogo.
    \item Identificar el funcionamiento del buscador.
    \item Revisar la informacion presentada para
    los productos.
    \item Explorar las opciones del carrito de compras.
    \item Identificar las opciones de seleccion
    de establecimientos.
    \item Registrar las evidencias visuales
    correspondientes.
\end{enumerate}

\subsection{Identificacion de necesidades}

A partir de las funcionalidades consideradas
se establecen las siguientes necesidades preliminares:

\begin{longtable}{p{1.5cm}p{3cm}p{6cm}p{3cm}}
\caption{Matriz de necesidades identificadas}\\
\toprule
\textbf{ID} &
\textbf{Actor} &
\textbf{Necesidad} &
\textbf{Fuente} \\
\midrule
\endfirsthead

\toprule
\textbf{ID} &
\textbf{Actor} &
\textbf{Necesidad} &
\textbf{Fuente} \\
\midrule
\endhead

NE-01 &
Cliente &
Consultar productos disponibles en el catalogo. &
Observacion del catalogo. \\

NE-02 &
Cliente &
Encontrar productos mediante una busqueda. &
Observacion del buscador. \\

NE-03 &
Cliente &
Consultar informacion de los productos. &
Observacion de productos. \\

NE-04 &
Cliente &
Agregar productos a una lista de compra. &
Observacion del carrito. \\

NE-05 &
Cliente &
Seleccionar un establecimiento para realizar
consultas relacionadas con la compra. &
Observacion de tiendas. \\

NE-06 &
Empresa &
Disponer de un canal digital para ofrecer
productos a los clientes. &
Contexto del sistema. \\

\bottomrule
\end{longtable}

\textbf{Alcance de los resultados:}
las necesidades anteriores son una interpretacion
de las funcionalidades identificadas en el alcance
de la practica. No representan entrevistas ni
declaraciones textuales de usuarios.

\subsection{Registro de evidencias}

Las evidencias deben demostrar las funcionalidades
observadas durante la exploracion.

\begin{table}[H]
\centering
\caption{Registro de evidencias de elicitacion}
\small
\begin{tabularx}{\textwidth}{p{2cm}p{4cm}Y}
\toprule
\textbf{Codigo} &
\textbf{Evidencia requerida} &
\textbf{Relacion con las necesidades} \\
\midrule

EV-01 &
Captura del catalogo de productos. &
NE-01 \\

EV-02 &
Captura del buscador. &
NE-02 \\

EV-03 &
Captura de la informacion de un producto. &
NE-03 \\

EV-04 &
Captura del carrito de compras. &
NE-04 \\

EV-05 &
Captura de la seleccion de tienda. &
NE-05 \\

\bottomrule
\end{tabularx}
\end{table}

% ==================================================
\section{Analisis de requisitos}
% ==================================================

\subsection{Analisis de ambiguedades}

Durante el analisis se consideran expresiones
generales que deben transformarse en requisitos
claros y verificables.

\begin{table}[H]
\centering
\caption{Analisis de ambiguedades}
\small
\begin{tabularx}{\textwidth}{p{4cm}YY}
\toprule
\textbf{Expresion inicial} &
\textbf{Problema identificado} &
\textbf{Formulacion mejorada} \\
\midrule

El sistema debe buscar rapidamente. &
Rapidamente no define un tiempo verificable. &
El tiempo de respuesta de la busqueda se medira
mediante pruebas bajo condiciones definidas. \\

El sistema debe mostrar buenos productos. &
Buenos no define una caracteristica objetiva. &
El sistema presentara los productos con la
informacion disponible en su catalogo. \\

El carrito debe funcionar correctamente. &
Correctamente es una expresion subjetiva. &
El sistema incorporara el producto seleccionado
al carrito y actualizara su cantidad. \\

\bottomrule
\end{tabularx}
\end{table}

\subsection{Analisis de conflictos}

Se identifican posibles conflictos que deben
considerarse durante el desarrollo y la validacion.

\begin{table}[H]
\centering
\caption{Matriz de conflictos}
\small
\begin{tabularx}{\textwidth}{p{3cm}YY}
\toprule
\textbf{Conflicto} &
\textbf{Descripcion} &
\textbf{Alternativa de resolucion} \\
\midrule

Disponibilidad &
El usuario espera encontrar productos disponibles,
pero la disponibilidad puede depender de la tienda. &
Considerar el establecimiento seleccionado al
consultar la disponibilidad. \\

Informacion del producto &
El usuario requiere informacion suficiente,
pero no todos los datos pueden estar disponibles. &
Definir los campos minimos de informacion
que deben presentarse. \\

Rendimiento &
El usuario espera respuestas rapidas, pero
pueden existir restricciones tecnicas. &
Establecer objetivos medibles y realizar
pruebas de rendimiento. \\

\bottomrule
\end{tabularx}
\end{table}

Los conflictos anteriores son situaciones potenciales
identificadas durante el analisis. No se afirma que
hayan sido reportados por personal de la empresa.

\subsection{Priorizacion mediante MoSCoW}

Se utiliza el modelo MoSCoW para clasificar
la importancia de los requisitos.

\begin{itemize}
    \item \textbf{Must:} requisito indispensable.
    \item \textbf{Should:} requisito importante.
    \item \textbf{Could:} requisito conveniente.
    \item \textbf{Won't:} requisito fuera del alcance
    de esta practica.
\end{itemize}

\begin{table}[H]
\centering
\caption{Priorizacion de requisitos}
\small
\begin{tabularx}{\textwidth}{p{2cm}p{7cm}p{3cm}}
\toprule
\textbf{ID} &
\textbf{Requisito} &
\textbf{Prioridad} \\
\midrule

RF-01 & Consultar el catalogo de productos. & Must \\

RF-02 & Buscar productos por nombre. & Must \\

RF-03 & Consultar informacion del producto. & Should \\

RF-04 & Agregar productos al carrito. & Must \\

RF-05 & Modificar cantidades del carrito. & Should \\

RF-06 & Seleccionar una tienda. & Should \\

RNF-01 & Mantener tiempos de respuesta medibles. & Should \\

RNF-02 & Presentar una interfaz comprensible. & Should \\

\bottomrule
\end{tabularx}
\end{table}

\subsection{Analisis de dependencias}

Las dependencias permiten determinar que
funcionalidades necesitan de otras para
ejecutarse correctamente.

\begin{table}[H]
\centering
\caption{Matriz de dependencias}
\small
\begin{tabularx}{\textwidth}{p{2cm}p{3cm}YY}
\toprule
\textbf{ID} &
\textbf{Requisito} &
\textbf{Dependencia} &
\textbf{Justificacion} \\
\midrule

RF-01 &
Catalogo &
Ninguna identificada &
Permite consultar los productos. \\

RF-02 &
Busqueda &
RF-01 &
La busqueda requiere informacion
de los productos. \\

RF-03 &
Informacion &
RF-01 &
La informacion corresponde a un
producto del catalogo. \\

RF-04 &
Carrito &
RF-03 &
El usuario selecciona un producto
para agregarlo. \\

RF-05 &
Cantidad &
RF-04 &
Se necesita un producto previamente
incorporado al carrito. \\

RF-06 &
Seleccion de tienda &
Ninguna identificada &
Permite establecer un contexto
de consulta. \\

\bottomrule
\end{tabularx}
\end{table}

% ==================================================
\section{Especificacion de requisitos}
% ==================================================

\subsection{Criterios de especificacion}

Los requisitos se documentan utilizando los
siguientes elementos:

\begin{itemize}
    \item Identificador unico.
    \item Nombre del requisito.
    \item Descripcion.
    \item Tipo de requisito.
    \item Prioridad.
    \item Fuente de informacion.
    \item Criterios de aceptacion.
\end{itemize}

Se consideran las recomendaciones de IEEE 29148
y la sintaxis EARS para mejorar la claridad
de los requisitos.

\subsection{Requisitos funcionales}

\begin{longtable}{p{1.5cm}p{3cm}p{7cm}p{2cm}}
\caption{Especificacion de requisitos funcionales}\\
\toprule
\textbf{ID} &
\textbf{Nombre} &
\textbf{Descripcion} &
\textbf{Prioridad} \\
\midrule
\endfirsthead

\toprule
\textbf{ID} &
\textbf{Nombre} &
\textbf{Descripcion} &
\textbf{Prioridad} \\
\midrule
\endhead

RF-01 &
Consultar catalogo &
El sistema debera permitir al cliente consultar
los productos presentados en el catalogo. &
Must \\

RF-02 &
Buscar productos &
Cuando el cliente introduzca un termino de busqueda,
el sistema debera presentar los productos que
coincidan con la consulta, segun la funcionalidad
disponible. &
Must \\

RF-03 &
Consultar informacion &
El sistema debera permitir al cliente visualizar
la informacion disponible de un producto
seleccionado. &
Should \\

RF-04 &
Agregar al carrito &
Cuando el cliente seleccione la opcion de agregar
un producto, el sistema debera incorporarlo
al carrito. &
Must \\

RF-05 &
Modificar cantidades &
El sistema debera permitir modificar la cantidad
de los productos incorporados al carrito,
cuando esta funcionalidad se encuentre disponible. &
Should \\

RF-06 &
Seleccionar tienda &
El sistema debera permitir seleccionar un
establecimiento mediante las opciones disponibles
en la plataforma. &
Should \\

\bottomrule
\end{longtable}

\subsection{Criterios de aceptacion}

Los siguientes criterios establecen condiciones
observables para determinar si los requisitos
funcionales se cumplen.

\begin{table}[H]
\centering
\caption{Criterios de aceptacion de requisitos}
\small
\begin{tabularx}{\textwidth}{p{2cm}YY}
\toprule
\textbf{ID} &
\textbf{Condicion} &
\textbf{Resultado esperado} \\
\midrule

RF-01 &
Dado que el cliente accede al catalogo. &
El sistema presenta los productos disponibles
en la interfaz. \\

RF-02 &
Dado que el cliente introduce un termino
de busqueda. &
El sistema presenta resultados relacionados
con el termino ingresado o informa que no
se encontraron coincidencias. \\

RF-03 &
Dado que el cliente selecciona un producto. &
El sistema presenta la informacion disponible
de dicho producto. \\

RF-04 &
Dado que el cliente selecciona un producto
y solicita agregarlo al carrito. &
El producto se incorpora al carrito y se
actualiza su contenido. \\

RF-05 &
Dado que existe un producto en el carrito. &
El sistema permite modificar su cantidad
y actualiza el contenido. \\

RF-06 &
Dado que el cliente accede a la seleccion
de establecimientos. &
El sistema presenta las opciones disponibles
para seleccionar una tienda. \\

\bottomrule
\end{tabularx}
\end{table}

\subsection{Requisitos no funcionales}

Los requisitos no funcionales establecen condiciones
de calidad que debe satisfacer el sistema.

\begin{longtable}{p{1.5cm}p{3cm}p{8cm}p{2cm}}
\caption{Especificacion de requisitos no funcionales}\\
\toprule
\textbf{ID} &
\textbf{Categoria} &
\textbf{Descripcion} &
\textbf{Prioridad} \\
\midrule
\endfirsthead

\toprule
\textbf{ID} &
\textbf{Categoria} &
\textbf{Descripcion} &
\textbf{Prioridad} \\
\midrule
\endhead

RNF-01 &
Rendimiento &
El sistema debera establecer un tiempo objetivo
de respuesta para las operaciones de consulta.
Este valor debera definirse y comprobarse
mediante pruebas de rendimiento. &
Should \\

RNF-02 &
Usabilidad &
La interfaz debera presentar etiquetas,
opciones y mensajes comprensibles para
facilitar la interaccion del usuario. &
Should \\

RNF-03 &
Compatibilidad &
La plataforma debera poder utilizarse desde
navegadores web compatibles y dispositivos
con diferentes resoluciones. &
Could \\

RNF-04 &
Seguridad &
La plataforma debera proteger la informacion
del usuario y controlar el acceso a las
funcionalidades que requieran autenticacion. &
Must \\

\bottomrule
\end{longtable}

Los requisitos no funcionales anteriores representan
objetivos propuestos para la especificacion.
Su cumplimiento no se considera demostrado mediante
la observacion visual de la interfaz.

\subsection{Trazabilidad de requisitos}

La trazabilidad permite relacionar cada requisito
con la necesidad que lo origina y con la evidencia
que permite justificar su identificacion.

\begin{table}[H]
\centering
\caption{Matriz de trazabilidad}
\small
\begin{tabularx}{\textwidth}{p{2cm}p{2cm}p{3cm}Y}
\toprule
\textbf{Requisito} &
\textbf{Necesidad} &
\textbf{Fuente} &
\textbf{Evidencia requerida} \\
\midrule

RF-01 &
NE-01 &
Observacion del catalogo &
EV-01 \\

RF-02 &
NE-02 &
Observacion del buscador &
EV-02 \\

RF-03 &
NE-03 &
Observacion del producto &
EV-03 \\

RF-04 &
NE-04 &
Observacion del carrito &
EV-04 \\

RF-05 &
NE-04 &
Exploracion del carrito &
EV-04 \\

RF-06 &
NE-05 &
Observacion de tiendas &
EV-05 \\

RNF-01 &
NE-01 &
Objetivo de rendimiento &
Prueba futura \\

RNF-02 &
NE-01 &
Objetivo de usabilidad &
Evaluacion futura \\

RNF-03 &
NE-01 &
Objetivo de compatibilidad &
Prueba futura \\

RNF-04 &
NE-06 &
Objetivo de seguridad &
Auditoria futura \\

\bottomrule
\end{tabularx}
\end{table}

% ==================================================
\section{Validacion de requisitos}
% ==================================================

\subsection{Criterios de validacion}

Para validar los requisitos se consideran
las siguientes propiedades:

\begin{itemize}

    \item \textbf{Claridad:} el requisito debe
    comprenderse sin interpretaciones ambiguas.

    \item \textbf{Completitud:} debe contener
    la informacion necesaria para su comprension.

    \item \textbf{Consistencia:} no debe contradecir
    otros requisitos.

    \item \textbf{Verificabilidad:} debe permitir
    determinar si se cumple mediante una revision
    o prueba.

    \item \textbf{Trazabilidad:} debe relacionarse
    con una necesidad o fuente identificada.

    \item \textbf{Factibilidad:} debe poder
    implementarse considerando las restricciones
    del sistema.

\end{itemize}

\subsection{Lista de verificacion}

Se utiliza la siguiente lista para revisar
la calidad documental de los requisitos.

\begin{table}[H]
\centering
\caption{Lista de verificacion de requisitos}
\small
\begin{tabularx}{\textwidth}{p{1.5cm}YYp{2cm}}
\toprule
\textbf{ID} &
\textbf{Criterio} &
\textbf{Comprobacion} &
\textbf{Resultado} \\
\midrule

V-01 &
Claridad &
Las descripciones utilizan expresiones
comprensibles y especificas. &
Cumple \\

V-02 &
Identificacion &
Los requisitos tienen identificadores unicos. &
Cumple \\

V-03 &
Priorizacion &
Los requisitos tienen una prioridad asignada. &
Cumple \\

V-04 &
Trazabilidad &
Los requisitos se relacionan con necesidades
y fuentes identificadas. &
Parcial \\

V-05 &
Aceptacion &
Los requisitos funcionales disponen de
criterios de aceptacion. &
Cumple \\

V-06 &
Verificabilidad &
Los requisitos permiten establecer mecanismos
para comprobar su cumplimiento. &
Parcial \\

V-07 &
Evidencias &
Las capturas originales respaldan cada
funcionalidad identificada. &
Pendiente \\

V-08 &
Factibilidad &
Los requisitos requieren una evaluacion
tecnica para confirmar su viabilidad. &
Pendiente \\

\bottomrule
\end{tabularx}
\end{table}

\textbf{Interpretacion de resultados:}

\begin{itemize}
    \item \textbf{Cumple:} el criterio se encuentra
    cubierto en la documentacion elaborada.

    \item \textbf{Parcial:} existe informacion
    relacionada, pero requiere comprobaciones
    adicionales.

    \item \textbf{Pendiente:} se necesita evidencia
    o una actividad adicional para confirmar
    el cumplimiento.
\end{itemize}

\subsection{Resultados de validacion}

La revision documental permite establecer que
los requisitos funcionales cuentan con identificadores,
descripciones y criterios de aceptacion.

La matriz de trazabilidad establece una relacion
entre las necesidades identificadas y los requisitos
propuestos.

Sin embargo, la validacion de las funcionalidades
frente al comportamiento real de la plataforma
requiere conservar y revisar las capturas obtenidas
durante la observacion.

De igual manera, los requisitos no funcionales
relacionados con rendimiento, seguridad y
compatibilidad requieren pruebas tecnicas
que no forman parte de la observacion visual.

Por tanto, los resultados corresponden a una
validacion documental preliminar y no a una
certificacion del funcionamiento interno
de la plataforma.

% ==================================================
\section{Diagrama del proceso de requisitos aplicado}
% ==================================================

\subsection{Representacion del proceso}

El siguiente diagrama representa la secuencia
de actividades aplicadas durante la practica,
incluyendo las actividades transversales de
gestion de requisitos y gestion de conflictos.

\begin{figure}[H]
\centering

% Desactivamos shorthands de babel para evitar conflictos con TikZ en flechas y colores
\shorthandoff{>}

\begin{tikzpicture}[
    node distance=1.25cm,
    proceso/.style={
        rectangle,
        rounded corners=4pt,
        draw=blue!60!black,
        fill=blue!8,
        text width=8cm,
        minimum height=1cm,
        align=center,
        font=\small
    },
    resultado/.style={
        rectangle,
        rounded corners=4pt,
        draw=green!50!black,
        fill=green!10,
        text width=8cm,
        minimum height=1cm,
        align=center,
        font=\small
    },
    transversal/.style={
        rectangle,
        rounded corners=4pt,
        draw=orange!70!black,
        fill=orange!12,
        text width=3.3cm,
        minimum height=1cm,
        align=center,
        font=\small
    },
    flecha/.style={
        -{Stealth[length=3mm]},
        thick,
        draw=black!70
    }
]

\node[proceso] (inicio) {
    \textbf{Inicio}\\
    Sistema seleccionado: Almacenes Tia
};

\node[proceso, below=of inicio] (elicitacion) {
    \textbf{1. Elicitacion}\\
    Observacion, exploracion e identificacion
    de necesidades
};

\node[proceso, below=of elicitacion] (analisis) {
    \textbf{2. Analisis}\\
    Ambiguedades, conflictos, prioridades
    y dependencias
};

\node[proceso, below=of analisis] (especificacion) {
    \textbf{3. Especificacion}\\
    Requisitos funcionales, no funcionales
    y criterios de aceptacion
};

\node[proceso, below=of especificacion] (validacion) {
    \textbf{4. Validacion}\\
    Lista de verificacion y revision
    de evidencias
};

\node[resultado, below=of validacion] (fin) {
    \textbf{Resultado}\\
    Requisitos documentados, trazables
    y revisados
};

\draw[flecha] (inicio) -- (elicitacion);
\draw[flecha] (elicitacion) -- (analisis);
\draw[flecha] (analisis) -- (especificacion);
\draw[flecha] (especificacion) -- (validacion);
\draw[flecha] (validacion) -- (fin);

\node[transversal, right=1.2cm of analisis] (gestion) {
    \textbf{Gestion de requisitos}\\
    Identificacion, trazabilidad
    y control de cambios
};

\node[transversal, right=1.2cm of especificacion] (conflictos) {
    \textbf{Gestion de conflictos}\\
    Identificacion y resolucion
    de discrepancias
};

\draw[<->, dashed, thick, orange!80!black]
    (gestion.west) -- (analisis.east);

\draw[<->, dashed, thick, orange!80!black]
    (gestion.south west) -- (especificacion.east);

\draw[<->, dashed, thick, orange!80!black]
    (conflictos.west) -- (especificacion.east);

\draw[<->, dashed, thick, orange!80!black]
    (conflictos.north west) -- (analisis.east);

\end{tikzpicture}

% Reactivamos los shorthands de español
\shorthandon{>}

\caption{Diagrama del proceso de ingenieria de
requisitos aplicado a Almacenes Tia.}
\label{fig:proceso-requisitos}
\end{figure}

\subsection{Descripcion de las fases}

\begin{enumerate}

    \item \textbf{Elicitacion:}

    Se explora la plataforma y se identifican
    las necesidades de los usuarios mediante
    la observacion de las funcionalidades.

    \item \textbf{Analisis:}

    Se examinan las necesidades identificadas,
    se detectan ambiguedades y se establecen
    prioridades y dependencias.

    \item \textbf{Especificacion:}

    Se documentan los requisitos funcionales
    y no funcionales, incorporando identificadores,
    fuentes y criterios de aceptacion.

    \item \textbf{Validacion:}

    Se revisa la calidad de los requisitos mediante
    criterios de claridad, consistencia,
    completitud y verificabilidad.

\end{enumerate}

Las actividades de gestion de requisitos y
gestion de conflictos se mantienen durante
las diferentes etapas del proceso.

% ==================================================
\section{Resultados y evidencias de la aplicacion}
% ==================================================

\subsection{Resultados obtenidos}

Como resultado de la aplicacion del proceso
de ingenieria de requisitos se elaboraron
los siguientes artefactos:

\begin{table}[H]
\centering
\caption{Matriz general de resultados}
\small
\begin{tabularx}{\textwidth}{p{3cm}YY}
\toprule
\textbf{Etapa} &
\textbf{Artefacto} &
\textbf{Resultado} \\
\midrule

Framework &
Matriz de actividades &
Relacion de actividades core y transversales
con el sistema. \\

Elicitacion &
Matriz de necesidades &
Identificacion de necesidades y fuentes. \\

Elicitacion &
Registro de evidencias &
Identificacion de las capturas requeridas. \\

Analisis &
Matriz de ambiguedades &
Reformulacion de expresiones generales. \\

Analisis &
Matriz MoSCoW &
Clasificacion de prioridades. \\

Analisis &
Matriz de dependencias &
Relacion entre requisitos. \\

Especificacion &
Matriz de requisitos &
Documentacion de requisitos funcionales
y no funcionales. \\

Especificacion &
Matriz de trazabilidad &
Relacion entre necesidades, requisitos y fuentes. \\

Validacion &
Lista de verificacion &
Revision documental de los requisitos. \\

Proceso &
Diagrama de actividades &
Representacion de la secuencia de trabajo. \\

\bottomrule
\end{tabularx}
\end{table}

\subsection{Evidencias fotograficas}

En este apartado se deben incorporar las capturas
reales obtenidas durante la exploracion del sitio web.

Las evidencias deben estar identificadas mediante
un codigo, titulo y descripcion.

% ==================================================
\section{Conclusiones}
% ==================================================

La aplicacion del framework de ingenieria de
requisitos de Klaus Pohl permitio organizar
el proceso de identificacion, analisis,
especificacion y validacion de requisitos
para la plataforma de comercio electronico
de Almacenes Tia.

Mediante la exploracion de la interfaz publica
se identificaron funcionalidades relacionadas
con el catalogo, la busqueda de productos,
la consulta de informacion, el carrito de compras
y la seleccion de establecimientos.

El analisis permitio organizar las necesidades,
identificar posibles ambiguedades, establecer
prioridades mediante el modelo MoSCoW y
determinar dependencias entre funcionalidades.

La especificacion permitio estructurar requisitos
funcionales y no funcionales, incorporando
identificadores, prioridades, fuentes y criterios
de aceptacion. Asimismo, la matriz de trazabilidad
facilita relacionar las necesidades con los
requisitos documentados.

Finalmente, se establecio una lista de verificacion
para revisar la calidad documental de los requisitos
y un diagrama que representa la secuencia del
proceso aplicado.

Se concluye que la ingenieria de requisitos
permite transformar las necesidades identificadas
en especificaciones organizadas y verificables,
facilitando la comprension del sistema y
proporcionando una base para futuras actividades
de desarrollo y evaluacion.

% ==================================================
\section{Referencias bibliograficas}
% ==================================================

\begin{thebibliography}{9}

\bibitem{pohl}
K. Pohl,
\textit{Requirements Engineering: Fundamentals,
Principles, and Techniques}.
Springer, 2010.

\bibitem{ieee}
IEEE,
\textit{ISO/IEC/IEEE 29148:2018:
Systems and Software Engineering ---
Life Cycle Processes --- Requirements Engineering}.
2018.

\bibitem{ears}
A. Mavin, P. Wilkinson, A. Harwood y M. Novak,
``Easy Approach to Requirements Syntax (EARS),''
en \textit{Proceedings of the 17th IEEE International
Requirements Engineering Conference}, 2009.

\bibitem{iso}
International Organization for Standardization,
\textit{ISO/IEC 25010:2011:
Systems and Software Engineering ---
Systems and Software Quality Requirements
and Evaluation (SQuaRE) --- System and Software
Quality Models}.
2011.

\bibitem{tia}
Almacenes Tia,
\textit{Plataforma oficial de comercio electronico}.
Disponible en:
\url{https://www.tia.com.ec/}

\end{thebibliography}

\end{document}